\begin{tikzpicture}[scale=0.7,
    dot/.style={circle,fill=black,inner sep=0.8pt},
    cross/.style={line width=0.7pt},
    every path/.style={line cap=round}
]

% Koordinaten (frei gewählt, nur für die Form)
\def\len{3.4}
\coordinate (A) at (-3,1);    % linker Kreuzpunkt
\coordinate (E) at (3,-1);     % rechter Kreuzpunkt
\coordinate (B) at ($(A) + (-35:\len)$);  % linker grüner Punkt
\coordinate (D) at ($(E) + (215:\len)$);  % rechter grüner Punkt
\coordinate (C) at ($(B)!0.5!(D)$);     % mittlerer grüner Punkt
\coordinate (Av) at (0,1.8); % senkrecht über C
\coordinate (Cv) at (0,-1.8); % senkrecht unter C

% definiere Farben
\definecolor{jointgreen}{RGB}{110,170,110}
\definecolor{jointyellow}{RGB}{255,204,102}

% orange Linien
\draw[jointyellow,very thick] (A) -- (B);
\draw[jointyellow,very thick] (D) -- (E);

% grüner Mittelsteg
\draw[jointgreen,very thick] (B) -- (C) -- (D);

% Punkte auf dem grünen Segment
\node[dot] at (B) {};
\node[dot] at (C) {};
\node[dot] at (D) {};

% senkrechte gestrichelte Linie unter dem mittleren Punkt
\draw[densely dotted] (Av) -- (Cv);

% kleine Kreuze an den Enden
\draw[cross] (A) ++(-0.06,0.06) -- ++(0.12,-0.12);
\draw[cross] (A) ++(-0.06,-0.06) -- ++(0.12,0.12);

\draw[cross] (E) ++(-0.06,0.06) -- ++(0.12,-0.12);
\draw[cross] (E) ++(-0.06,-0.06) -- ++(0.12,0.12);

\end{tikzpicture}
